%-------------------------
% Resume in Latex
% Author
% License : MIT
%------------------------

%---- Required Packages and Functions ----

\documentclass[a4paper,11pt]{article}
\usepackage{latexsym}
\usepackage{xcolor}
\usepackage{float}
\usepackage{ragged2e}
\usepackage[empty]{fullpage}
\usepackage{wrapfig}
\usepackage{lipsum}
\usepackage{tabularx}
\usepackage{titlesec}
\usepackage{geometry}
\usepackage{marvosym}
\usepackage{verbatim}
\usepackage{enumitem}
\usepackage[hidelinks]{hyperref}
\usepackage{fancyhdr}
\usepackage{fontawesome5}
\usepackage{multicol}
\usepackage{graphicx}
\usepackage{cfr-lm}
\usepackage[T1]{fontenc}
\setlength{\multicolsep}{0pt} 
\pagestyle{fancy}
\fancyhf{} % clear all header and footer fields
\fancyfoot{}
\renewcommand{\headrulewidth}{0pt}
\renewcommand{\footrulewidth}{0pt}
\geometry{left=1.4cm, top=0.8cm, right=1.2cm, bottom=1cm}
% Adjust margins
%\addtolength{\oddsidemargin}{-0.5in}
%\addtolength{\evensidemargin}{-0.5in}
%\addtolength{\textwidth}{1in}
\usepackage[most]{tcolorbox}
\tcbset{
	frame code={}
	center title,
	left=0pt,
	right=0pt,
	top=0pt,
	bottom=0pt,
	colback=gray!20,
	colframe=white,
	width=\dimexpr\textwidth\relax,
	enlarge left by=-2mm,
	boxsep=4pt,
	arc=0pt,outer arc=0pt,
}

\urlstyle{same}

\raggedright{}
\setlength{\tabcolsep}{0in}
% Fancyhdr settings
\setlength{\footskip}{4.5pt}



% Sections formatting
\titleformat{\section}{
  \vspace{-4pt}\scshape\raggedright\large
}{}{0em}{}[\color{black}\titlerule{} \vspace{-7pt}]

%-------------------------
% Custom commands
\newcommand{\resumeItem}[2]{
  \item{
    \textbf{#1}{\hspace{0.5mm}#2 \vspace{-0.5mm}}
  }
}

\newcommand{\resumePOR}[3]{
\vspace{0.5mm}\item
    \begin{tabular*}{0.97\textwidth}[t]{l@{\extracolsep{\fill}}r}
        \textbf{#1}\hspace{0.3mm}#2 & \textit{\small{#3}} 
    \end{tabular*}
    \vspace{-2mm}
}

\newcommand{\resumeSubheading}[4]{
\vspace{0.5mm}\item
    \begin{tabular*}{0.98\textwidth}[t]{l@{\extracolsep{\fill}}r}
        \textbf{#1} & \textit{\footnotesize{#4}} \\
        \textit{\footnotesize{#3}} &  \footnotesize{#2}\\
    \end{tabular*}
    \vspace{-2.4mm}
}

\newcommand{\resumeProject}[4]{
\vspace{0.5mm}\item
    \begin{tabular*}{0.98\textwidth}[t]{l@{\extracolsep{\fill}}r}
        \textbf{#1} & \textit{\footnotesize{#3}} \\
        \footnotesize{\textit{#2}} & \footnotesize{#4}
    \end{tabular*}
    \vspace{-2.4mm}
    \par
}

\newcommand{\resumeSubItem}[2]{\resumeItem{#1}{#2}\vspace{-4pt}}
% \renewcommand{\labelitemii}{$\circ$}
\renewcommand{\labelitemi}{$\vcenter{\hbox{\tiny$\bullet$}}$}
\newcommand{\resumeSubHeadingListStart}{\begin{itemize}[leftmargin=*,labelsep=0mm]}
\newcommand{\resumeHeadingSkillStart}{\begin{itemize}[leftmargin=*,itemsep=1.7mm, rightmargin=2ex]}
\newcommand{\resumeItemListStart}{\begin{justify}\begin{itemize}[leftmargin=3ex, rightmargin=2ex, noitemsep,labelsep=1.2mm,itemsep=0mm]\small}
\newcommand{\resumeSubHeadingListEnd}{\end{itemize}\vspace{2mm}}
\newcommand{\resumeHeadingSkillEnd}{\end{itemize}\vspace{-2mm}}
\newcommand{\resumeItemListEnd}{\end{itemize}\end{justify}\vspace{-2mm}}
\newcommand{\cvsection}[1]{%
\vspace{2mm}
\begin{tcolorbox}
    \textbf{\large #1}
\end{tcolorbox}
    \vspace{-4mm}
}
\newcolumntype{L}{>{\raggedright\arraybackslash}X}%
\newcolumntype{R}{>{\raggedleft\arraybackslash}X}%
\newcolumntype{C}{>{\centering\arraybackslash}X}%
%---- End of Packages and Functions ------

%-------------------------------------------
%%%%%%  CV STARTS HERE  %%%%%%%%%%%
%%%%%% DEFINE ELEMENTS HERE %%%%%%%
\newcommand{\name}{Agus Fuad Mudhofar} % Your Name
\newcommand{\course}{Teknik Komputer} % Your Program
% \newcommand{\roll}{xxxxxxx} % Your Roll No.
\newcommand{\phone}{8882111890} % Your Phone Number
\newcommand{\emaila}{agusfuad090@gmail.com} %Email 1

\begin{document}
\fontfamily{cmr}\selectfont
%----------HEADING-----------------

{
  \begin{tabularx}{\linewidth}{L r}                                                                                                                                                    \\
    \textbf{\Large \name}                           & {\raisebox{0.0\height}{\footnotesize \faPhone}\ +62-\phone}                                                        \\
    {Institut Teknologi Sepuluh Nopember, Surabaya} & \href{mailto:\emaila}{\raisebox{0.0\height}{\footnotesize \faEnvelope}\ {\emaila}}                                 \\
                                                    & \href{https://maedakenji.github.io/Portofolio/}{\raisebox{0.0\height}{\footnotesize \faGithub}\ {Profil GitHub}}             \\
                                                    & \href{https://www.linkedin.com/in/agus4434/}{\raisebox{0.0\height}{\footnotesize \faLinkedin}\ {Profil LinkedIn}}
  \end{tabularx}
}

%-----------EDUCATION-----------
\section{\textbf{Pendidikan}}
\resumeSubHeadingListStart
\resumeSubheading
{Sarjana Teknik Komputer}{IPK: 3.64}
{Institut Teknologi Sepuluh Nopember, Surabaya}{2022 -- Sekarang}
\resumeSubHeadingListEnd
\vspace{-5.5mm}
%

%-----------PROJECTS-----------------

\section{\textbf{Proyek Pribadi}}
\resumeSubHeadingListStart

\resumeProject
{Sappy: Aplikasi Manajemen Peternakan Multi-Peran (Flutter)} % Project Name
{Aplikasi manajemen peternakan untuk berbagai peran (Peternak, Admin, Dokter) dengan dukungan NFC.}
{} % Event Dates
{}
\resumeItemListStart
\item {Menyediakan akses berbasis peran dan dasbor untuk peternak, administrator, dan dokter, semuanya dalam satu aplikasi seluler terpadu.}
\item {Mendukung bahasa Inggris dan Indonesia menggunakan lokalisasi bawaan.}
\item {Mengimplementasikan manajemen state melalui Provider dan mengaktifkan fungsionalitas NFC untuk identifikasi dan pelacakan yang aman.}
\item {Mencakup visualisasi data, autentikasi, dan konfigurasi lingkungan menggunakan paket ekosistem Flutter modern.}
\item {Teknologi yang Digunakan: Flutter, Dart, Provider, NFC, fl\_chart, carousel\_slider, font/aset kustom.}
\resumeItemListEnd
\vspace{-2mm}


\resumeProject
{Aplikasi Web Membaca Komik} % Project Name
{Platform membaca komik berbasis web yang mendukung manga, manhua, dan manhwa.} % Project Description
{} % Event Dates
{}
\resumeItemListStart
\item {Mengembangkan dua arsitektur full-stack terpisah (Laravel-React dan Express-React) untuk melakukan tolok ukur (benchmark) kinerja framework dalam lingkungan rendering gambar dengan lalu lintas tinggi.}
\item {Mengimplementasikan fitur membaca inti termasuk manajemen chapter dinamis, pelacakan riwayat membaca, autentikasi, dan tata letak responsif untuk berbagai perangkat.}
\item {Mengintegrasikan MinIO sebagai solusi penyimpanan objek lokal yang kompatibel dengan S3 untuk penyimpanan mandiri (self-hosted) halaman komik berukuran besar secara cepat dan efisien.}
\item {Mendeploy iterasi Laravel menggunakan FrankenPHP sebagai server aplikasi modern untuk memaksimalkan kecepatan penanganan permintaan (request) dan mengoptimalkan runtime backend.}
\item {Teknologi yang Digunakan: React, Laravel, Express.js, PostgreSQL, FrankenPHP, MinIO, dengan arsitektur komponen modular dan penanganan aset otomatis.}
\resumeItemListEnd
\vspace{-2mm}

% End of Projects
\resumeSubHeadingListEnd
\vspace{-4mm}

% %-----------EXPERIENCE-----------------

\section{\textbf{Pengalaman}}
\resumeSubHeadingListStart


\resumeSubheading
{Magang Riset -- Pengembang Mobile \& Computer Vision Engineer}{Badan Riset dan Inovasi Nasional (BRIN)}
{Bandung}{Feb -- Jul 2026}

\vspace{0mm}
\begin{itemize}[leftmargin=1.5ex, label={}]

\item{a. Pengembangan Aplikasi Full-Stack}
    \begin{itemize}[label=-]
        \item {Membangun aplikasi Flutter lintas platform (Mobile, Tablet, Desktop) untuk menghitung kebutuhan nutrisi NPK tanaman sayuran, dideploy ke berbagai ukuran layar.}
        \item {Membangun backend RESTful API dengan Laravel 12 dan PostgreSQL; mengimplementasikan autentikasi berbasis token yang aman menggunakan Laravel Sanctum.}
        \item {Merancang UI yang intuitif dan responsif di Flutter, memastikan pengalaman pengguna yang mulus di semua faktor bentuk perangkat.}
    \end{itemize}

\item{b. Sistem AI \& Computer Vision}
    \begin{itemize}[label=-]
        \item {Mengembangkan sistem visual awareness menggunakan input kamera dan model deep learning untuk mengekstrak informasi lingkungan yang kontekstual, termasuk atribut pengguna (misalnya, pakaian) dan objek sekitar.}
        \item {Mengintegrasikan output computer vision dengan Large Language Model (LLM) untuk menghasilkan pemahaman kontekstual yang diperkaya, memungkinkan interpretasi aktivitas pengguna seperti berjalan, bersepeda, berkendara motor, dan mengemudi.}
        \item {Merancang dan mengembangkan modul khusus untuk pipeline pemrosesan, mengintegrasikan deteksi objek, pengenalan aksi manusia, dan semantic reasoning untuk meningkatkan situational awareness.}
    \end{itemize}
\end{itemize}

\newpage
\resumeSubheading
{Asisten Praktikum Lab -- Pemrograman Dasar}{Multimedia and IoT Laboratory (MIOT)}
{Surabaya}{Agu 2025 -- Nov 2025}
\resumeItemListStart
\item {Mengembangkan dan memelihara materi lab, termasuk latihan pengodean dan panduan pemecahan masalah selama praktikum, serta bertanggung jawab untuk mendeploy dan memelihara situs web online judge DMOJ yang digunakan dalam perkuliahan.}
\item {Membimbing mahasiswa melalui tugas-tugas praktikum dalam C/C++, memberikan dukungan real-time dan penjelasan teknis.}
\item {Mengevaluasi kiriman tugas mahasiswa, memberikan umpan balik konstruktif, dan berkontribusi pada rubrik penilaian untuk memastikan evaluasi yang adil dan konsisten.}
\item {Mengorganisasi dan memimpin sesi tinjauan sebelum ujian, mengatasi kesalahpahaman umum dan memperkuat konsep pemrograman utama.}
\resumeItemListEnd
\vspace{-2mm}


\resumeSubheading
{Pengembang Web -- Halaman Web MIOT}
{Multimedia and IoT Laboratory (MIOT)}
{Surabaya}{Agu 2024 -- Agu 2025}
\resumeItemListStart
\item {Merancang dan mengembangkan halaman web resmi laboratorium MIOT untuk menampilkan riset, publikasi, dan anggota tim.}
\item {Mengimplementasikan antarmuka pengguna yang responsif dan modern menggunakan React.js dan Tailwind CSS, memastikan aksesibilitas di berbagai perangkat.}
\item {Mengintegrasikan manajemen konten dinamis untuk berita, acara, dan profil anggota menggunakan struktur data berbasis JSON.}
\item {Berkolaborasi dengan anggota lab untuk mengumpulkan kebutuhan, melakukan iterasi pada desain, dan mendeploy situs web menggunakan GitHub Pages.}
\item {Mengoptimalkan performa situs web dan SEO, serta menyediakan dokumentasi untuk pemeliharaan dan pembaruan di masa mendatang.}
\resumeItemListEnd
\vspace{-4mm}




% End of Experience
\resumeSubHeadingListEnd
\vspace{-5.5mm}

\section{\textbf{Kepanitiaan}}
\resumeSubHeadingListStart


\resumeSubheading
{Ketua Divisi Logistik -- Multimedia and Game Event (MAGE X)}
{Multimedia and Game Event (MAGE X)}
{Surabaya}{Jan -- Nov 2024}
\resumeItemListStart
\item {Memimpin tim logistik untuk MAGE X, sebuah acara kompetisi multimedia dan game tingkat universitas yang besar.}
\item {Mengkoordinasikan pengadaan, manajemen inventaris, dan distribusi materi serta peralatan acara.}
\item {Menyusun dan mengeksekusi rencana logistik untuk memastikan penyiapan yang tepat waktu dan kelancaran operasi di berbagai lokasi acara.}
\item {Berkolaborasi dengan departemen lain untuk memenuhi kebutuhan real-time dan menyelesaikan tantangan logistik selama acara.}
\item {Mengelola hubungan dengan vendor, pelacakan anggaran, dan rekonsiliasi inventaris setelah acara.}
\resumeItemListEnd


\resumeSubheading
{PIC IoT Division – TIM Kawal GEMASTIK}{Institut Teknologi Sepuluh Nopember (ITS)}
{Surabaya}{Jan 2025 \- Nov 2025}
\resumeItemListStart
\item {Bertindak sebagai koordinator antara dosen kawal dengan tim mahasiswa.}
\item {Membantu proses seleksi tim melalui evaluasi proposal, penilaian teknis, dan kesiapan untuk kompetisi.}
\item {Memberikan dukungan administratif untuk tim finalis, seperti pengecekan kesesuaian berkas, persiapan berkas pendaftaran, dan koordinasi dengan panitia nasional.}
\item {Mengkoordinasikan kegiatan mentoring dengan menjadwalkan, mengorganisasi, dan memfasilitasi sesi mentoring bersama para dosen kawal di bidang IoT.}
\item {Mendukung kebutuhan non\-teknis tim seperti mencarikan tim desain presentasi, pendampingan untuk 3D printing, dan pendampingan membuat video presentasi.}
\resumeItemListEnd


% End of Positions of Responsibility
\resumeSubHeadingListEnd

\section{\textbf{Sertifikasi}}
\resumeSubHeadingListStart
\resumeSubheading
{Sertifikasi Kefasihan Bahasa Inggris -- CEFR Level B2}{}
{DynEd International}{2024}
\resumeItemListStart
\item {Telah mencapai kefasihan bahasa Inggris setara dengan Common European Framework of Reference for Languages (CEFR) Level B2.}
\resumeItemListEnd
\resumeSubHeadingListEnd
\vspace{-16pt}


% %-----------SKILLS & INTERESTS-----------------

\section{\textbf{Keahlian Teknis dan Minat}}
\begin{itemize}[leftmargin=0.05in, label={}]
  \small{\item{
        \textbf{Bahasa Pemrograman}{: C/C++, Python, Javascript, HTML+CSS } \\
        \textbf{Pustaka / Library}{: C++ STL, Library Python, ReactJs }\\
        \textbf{Alat Pengembangan Web}{: Nodejs, VS Code, Git, GitHub } \\
        \textbf{Framework}{: ReactJs } \\
        \textbf{Cloud/Basis Data}{: MongoDb, Firebase, Basis Data Relasional (MySQL) } \\

        \textbf{Mata Kuliah Relevan}{: Struktur Data \& Algoritma, Sistem Operasi, Pemrograman Berorientasi Objek (OOP), Sistem Manajemen Basis Data, Rekayasa Perangkat Lunak } \\
        \textbf{Minat}{: Pengembangan Perangkat Lunak} \\
        \textbf{Soft Skills}{: Pemecahan Masalah, Pembelajaran Mandiri (Self-learning), Presentasi, Adaptabilitas} \\
        \textbf{Bahasa}{: Indonesia (Native), Inggris (B2 CEFR - Bersertifikat DynEd)} \\
        }}
\end{itemize}
\vspace{-16pt}


\end{document}
